\def\thoigian{2}
\section*
{Luyện tập chung}

\subsection{Mục tiêu}

\paragraph{Về kiến thức, kĩ năng}

\begin{itemize}
	\item Luyện tập giải hệ phương trình bậc nhất hai ẩn bằng phương pháp thế hay phương pháp cộng đại số.
	\item Luyện tập tìm nghiệm của hệ hai phương trình bậc nhất hai ẩn bằng máy tính cầm tay.
\end{itemize}

\paragraph{Về năng lực}

\begin{itemize}
	\item Rèn luyện năng lực toán học, nói riêng là năng lực tư duy và lập luận toán học, năng lực sử dụng công cụ, phương tiện học toán.
	\item Bồi dưỡng hứng thú học tập, ý thức làm việc nhóm, ý thức tìm tòi, khám phá và sáng tạo cho HS. 
\end{itemize}

\paragraph{Về phẩm chất}

\begin{itemize}
	\item Góp phần giúp HS rèn luyện và phát triển các phẩm chất tốt đẹp (yêu nước, nhân ái, chăm chỉ, trung thực, trách nhiệm):
	\begin{itemize}
		\item Tích cực phát biểu, xây dựng bài và tham gia các hoạt động nhóm;
		\item Có ý thức tích cực tìm tòi, sáng tạo trong học tập; phát huy điểm mạnh, khắc phục các điểm yếu của bản thân. 
	\end{itemize}
\end{itemize}

\subsection{Thiết bị dạy học và học liệu}

\paragraph{Giáo viên}

\begin{itemize}
	\item Giáo án, bảng phụ, máy chiếu (nếu có), phiếu học tập, … 
\end{itemize}

\paragraph{Học sinh}

\begin{itemize}
	\item SGK, vở ghi, dụng cụ học tập, máy tính cầm tay. 
\end{itemize}

\subsection{Tiến trình dạy học}

Bài học này dạy trong 02 tiết:

\begin{itemize}
	\item Tiết 1: Ôn lại lí thuyết và các ví dụ.
	\item Tiết 2: Các bài tập cuối bài.
\end{itemize}

\subsubsection{Ôn tập lý thuyết và các ví dụ}

\paragraph{Hoạt động khởi động}

\subparagraph{Mục tiêu}

\begin{itemize}
	\item Hệ thống lại kiến thức về các phương pháp giải hệ hai phương trình bậc nhất hai ẩn.
\end{itemize}

\begin{paracol}{2}
	\subparagraph{Nội dung}
	\switchcolumn
	\subparagraph{Sản phẩm}
\end{paracol}

\begin{ontap}
	Nhắc lại các bước để giải hệ hai phương trình bậc nhất hai ẩn bằng:
	\begin{listEX}
		\item phương pháp thế;
		\item phương pháp cộng đại số.
	\end{listEX}
	%%%
	\doicot
	%%%
	\begin{listEX}
		\item Giải hệ phương trình bằng phương pháp thế:
		\begin{itemize}
			\item \textit{Bước 1.} Từ một phương trình của hệ, biểu diễn một ẩn theo ẩn kia rồi thế vào phương trình còn lại của hệ để được phương trình chỉ còn chứa một ẩn.
			\item \textit{Bước 2.} Giải phương trình một ẩn vừa nhận được, từ đó suy ra nghiệm của hệ đã cho.
		\end{itemize}
		\item Giải hệ phương trình bằng phương pháp cộng đại số:
		\begin{itemize}
			\item \textit{Bước 1.} Cộng hay trừ từng vế của hai phương trình trong hệ để được phương trình chỉ còn chứa một ẩn.
			\item \textit{Bước 2.} Giải phương trình một ẩn vừa nhận được, từ đó suy ra nghiệm của hệ phương trình đã cho.
		\end{itemize}
	\end{listEX}
\end{ontap}

\subparagraph{Tổ chức thực hiện}

\begin{tcth}
	\buocmot
	\item GV cho HS nhắc lại các bước để giải hệ hai phương trình bậc nhất hai ẩn bằng phương pháp thế và phương pháp cộng đại số.
	\buochai
	\item HS suy nghĩ, chuẩn bị câu trả lời.
	\buocba
	\item HS giơ tay phát biểu.
	\item Các HS khác nhận xét.
	\buocbon
	\item GV nhận xét và nhắc lại các lưu ý khi sử dụng các phương pháp này.
\end{tcth}

\paragraph{Hoạt động luyện tập}

\subparagraph{Mục tiêu}

\begin{itemize}
	\item Củng cố kĩ năng giải hệ hai phương trình bậc nhất hai ẩn.
\end{itemize}

\begin{paracol}{2}
	\subparagraph{Nội dung}
	\switchcolumn
	\subparagraph{Sản phẩm}
\end{paracol}

\begin{vidu}
	Giải hệ phương trình $\heva{&0,5x + 0,6y = 0,4 \\ &0,4x - 0,9y = 1,7.}$
	%%%
	\doicot
	%%%
	Nhân hai vế của mỗi phương trình với 10, ta được:
	$$\heva{&5x + 6y = 4 \\ &4x - 9y = 17.} \quad (1)$$
	Ta giải hệ (1). Nhân hai vế của phương trình thứ nhất với 3 và nhân hai vế của phương trình thứ hai với 2, ta được hệ:
	$$\heva{&15x + 18y = 12 \\ &8x - 18y = 34.} \quad (2)$$
	Cộng từng vế hai phương trình của hệ (2) ta được $23x = 46$, suy ra $x = 2$.\\
	Thế $x = 2$ vào phương trình thứ nhất của hệ (1), ta được $5 \cdot 2 + 6y = 4$ hay $6y = -6$, suy ra $y = -1$.\\
	Vậy hệ phương trình đã cho có nghiệm là $(2; -1)$.
\end{vidu}

\begin{vidu}
	Tìm các hệ số $x, y$ trong phản ứng hoá học đã được cân bằng sau:
	$$\ce{3Fe + xO2 -> yFe3O4}.$$
	%%%
	\doicot
	%%%	
	Vì số nguyên tử của $\text{Fe}$ và $\text{O}$ ở cả hai vế của phương trình phản ứng phải bằng nhau nên ta có hệ phương trình $\heva{&3 = 3y \\ &2x = 4y}$ hay $\heva{&y = 1 \\ &x = 2y.}$\\
	Giải hệ này ta được $x = 2, y = 1$.
\end{vidu}

\begin{vidu}
	Tìm hai số $a$ và $b$ để đường thẳng $y = ax + b$ đi qua hai điểm $A(-2; -1)$ và $B(2; 3)$.
	%%%
	\doicot
	%%%
	Đường thẳng $y = ax + b$ đi qua điểm $A(-2; -1)$ nên $a(-2) + b = -1$ hay $-2a + b = -1$.\\
	Tương tự, đường thẳng $y = ax + b$ đi qua điểm $B(2; 3)$ nên $a \cdot 2 + b = 3$ hay $2a + b = 3$.\\
	Từ đó, ta có hệ phương trình với hai ẩn là $a$ và $b$:
	$$\heva{&-2a + b = -1 \\ &2a + b = 3.}$$
	Cộng từng vế hai phương trình của hệ, ta được $2b = 2$, suy ra $b = 1$.\\
	Thay $b = 1$ vào phương trình thứ nhất, ta có $-2a + 1 = -1$, suy ra $a = 1$.\\
	Vậy với $a = 1, b = 1$ thì đường thẳng $y = ax + b$ đi qua hai điểm $A, B$ đã cho.
\end{vidu}

\subparagraph{Tổ chức thực hiện}

\noindent\faBookmarkO~\textit{\textbf{Ví dụ 1:}}

\begin{tcth}
	\buocmot
	\item GV yêu cầu HS làm việc cá nhân.
	\buochai
	\item HS thực hiện theo hướng dẫn của GV.
	\buocba
	\item Sau đó, GV gọi HS	lên bảng trình bày lời giải.
	\item Các HS khác nhận xét.
	\buocbon
	\item Sau đó GV tổng kết và lưu ý sai lầm thường mắc cho HS.
\end{tcth}

\noindent\faBookmarkO~\textit{\textbf{Ví dụ 2:}}

\begin{tcth}
	\buocmot
	\item GV tổ chức để HS nhớ lại các kiến thức Hóa học về cân bằng phương trình.
	\item GV yêu cầu HS làm việc nhóm đôi.
	\buochai
	\item HS thực hiện theo hướng dẫn của GV.
	\buocba
	\item Sau đó, GV gọi HS	lên bảng trình bày lời giải.
	\item Các HS khác nhận xét.
	\buocbon
	\item Sau đó GV tổng kết và lưu ý sai lầm thường mắc cho HS.
\end{tcth}

\noindent\faBookmarkO~\textit{\textbf{Ví dụ 3:}}

\begin{tcth}
	\buocmot
	\item GV tổ chức để HS nhớ lại các kiến thức Hình học về đồ thị hàm số bậc nhất, GV cần yêu cầu HS nhắc lại	về điều kiện để một điểm cho trước thuộc đồ thị hàm số bậc nhất.
	\item GV yêu cầu HS làm việc cá nhân.
	\buochai
	\item HS thực hiện theo hướng dẫn của GV.
	\buocba
	\item Sau đó, GV gọi HS	lên bảng trình bày lời giải.
	\item Các HS khác nhận xét.
	\buocbon
	\item Sau đó GV tổng kết và lưu ý sai lầm thường mắc cho HS.
\end{tcth}

\paragraph{Tổng kết và hướng dẫn công việc về nhà}

GV tổng kết lại nội dung bài học và dặn dò công việc ở nhà cho HS:

\begin{itemize}
	\item GV tổng kết lại các kiến thức trọng tâm của bài học: Các cách giải hệ hai phương trình bậc nhất hai ẩn.
	\item Nhắc HS về nhà ôn tập các nội dung đã học.
\end{itemize}

\subsubsection{Các bài tập cuối bài}

\paragraph{Hoạt động luyện tập}

\subparagraph{Mục tiêu}

\begin{itemize}
	\item Củng cố giải hệ phương trình bậc nhất hai ẩn bằng phương pháp thế, phương pháp cộng đại số hoặc sử dụng máy tính cầm tay.
\end{itemize}

\begin{paracol}{2}
	\subparagraph{Nội dung}
	\switchcolumn
	\subparagraph{Sản phẩm}
\end{paracol}

\begin{baitap}[1.10]
	Cho hai phương trình:
	\begin{align*}
		-2x + 5y &= 7; \quad (1) \\
		4x - 3y &= 7. \quad (2)
	\end{align*}
	Trong các cặp số $(2; 0)$, $(1; -1)$, $(-1; 1)$, $(-1; 6)$, $(4; 3)$ và $(-2; -5)$, cặp số nào là:
	\begin{listEX}
		\item Nghiệm của phương trình (1)?
		\item Nghiệm của phương trình (2)?
		\item Nghiệm của hệ gồm phương trình (1) và \\
		phương trình (2)?
	\end{listEX}
	%%%
	\doicot
	%%%
	\huongdan
	\begin{listEX}
		\item Nghiệm của phương trình (1): $(-1 ; 1)$, $(4 ; 3)$.
		\item Nghiệm của phương trình (2): $(1 ; -1)$, $(4 ; 3)$, $(- 2; -5)$.
		\item Nghiệm của hệ gồm (1) và (2): $(4 ; 3)$.
	\end{listEX}
\end{baitap}

\begin{baitap}[1.11]
	Giải các hệ phương trình sau bằng phương pháp thế:
	\begin{listEX}
		\item $\heva{&2x - y = 1 \\ &x - 2y = -1;}$
		\item $\heva{&0,5x - 0,5y = 0,5 \\ &1,2x - 1,2y = 1,2;}$
		\item $\heva{&x + 3y = -2 \\ &5x - 4y = 28.}$
	\end{listEX}
	%%%
	\doicot
	%%%
	\huongdan
	\begin{listEX}
		\item $(1;1)$;
		\item Vô số nghiệm: $(x;x-1)$ với $x\in\mathbb{R}$;
		\item $(4;-2)$.
	\end{listEX}
\end{baitap}

\begin{baitap}[1.12]
	Giải các hệ phương trình sau bằng phương pháp cộng đại số:
	\begin{listEX}
		\item $\heva{&5x + 7y = -1 \\ &3x + 2y = -5;}$
		\item $\heva{&2x - 3y = 11 \\ &-0,8x + 1,2y = 1;}$
		\item $\heva{&4x - 3y = 6 \\ &0,4x + 0,2y = 0,8.}$
	\end{listEX}
	%%%
	\doicot
	%%%
	\huongdan
	\begin{listEX}
		\item $(-3;2)$;
		\item Vô nghiệm;
		\item $(1{,}8;0{,}4)$.
	\end{listEX}
\end{baitap}

\begin{baitap}[1.13]
	Tìm các hệ số $x$, $y$ trong phản ứng hoá học đã được cân bằng sau:
	$$\ce{4Al + xO2 -> yAl2O3}.$$
	%%%
	\doicot
	%%%
	\huongdan\\
	$x=3$, $y=2$.
\end{baitap}

\begin{baitap}[1.14]
	Tìm $a$ và $b$ sao cho hệ phương trình 
	$$\heva{&ax + by = 1 \\ &ax + (2 - b)y = 3}$$ 
	có nghiệm là $(1; -2)$.
	%%%
	\doicot
	%%%
	\huongdan\\
	Vì $(x; y) = (1; -2)$ là nghiệm của hệ đã cho nên ta có:
	$$\heva{&a - 2b = 1 \\ &a + 2(b - 2) = 3} \quad \text{hay} \quad \heva{&a - 2b = 1 \\ &a + 2b = 7.}$$
	Giải hệ này ta được $a = 4, b = \dfrac{3}{2}$.
\end{baitap}

\subparagraph{Tổ chức thực hiện}

\textit{(Dùng chung cho các bài tập từ 1.10 đến 1.14)}

\begin{tcth}
	\buocmot
	\item GV tổ chức cho HS làm cá nhân.
	\buochai
	\item HS làm bài vào vở
	\item GV mời HS lên bảng làm bài.
	\buocba
	\item Các HS khác nhận xét bài làm của bạn.
	\buocbon
	\item GV nhận xét, chỉnh sửa nếu cần. 
	\item HS sửa bài.
\end{tcth}

\paragraph{Tổng kết và hướng dẫn cộng việc ở nhà}

GV tổng kết lại nội dung bài học và dặn dò công việc ở nhà cho HS:

\begin{itemize}
	\item GV tổng kết lại các kiến thức trọng tâm của bài học: Tóm tắt cách giải hệ hai phương trình bậc nhất hai ẩn bằng phương pháp thế, phương pháp cộng đại số, dùng máy tính cầm tay.
\end{itemize}